\begin{center}
{\fontsize{18}{24}\selectfont\bfseries 考虑变物性与径向收缩的药材\par 热湿传递及烘干时长预测}\par
\vspace{0.7em}
{\large\bfseries 摘\quad 要}
\end{center}
\noindent\hspace*{\parindent}%
针对圆柱形药材热风烘干过程中温湿分布不均、物性变化与径向收缩问题，本文建立由固定域向收缩域递进的热湿传递模型，预测径向温湿分布及全域达标时刻。采用保形分段三次Hermite插值（PCHIP）构造环境与半径的时间输入，以\textbf{表面加密有限体积法}离散守恒方程，并用\textbf{隐式变阶后向差分法（BDF）}求解。

\textbf{针对问题一，}建立常热物性导热方程与含水率相关的\textbf{非线性Fick扩散方程}，施加中心零通量和表面对流交换边界。1800 s时，中心、表面温度分别为33.5758、36.7879 $^\circ$C，干基含水率分别为2.5500、1.5102 kg/kg，表明预热期外层先升温、先失水，中心失水滞后。

\textbf{针对问题二，}引入含水率相关的热物性及温度、含水率共同决定的扩散系数，联合求解\textbf{双向耦合}的温度场与水分场。3 h时中心温度为49.8494 $^\circ$C，干基含水率仍为1.7662 kg/kg，说明温度接近稳定不能代表干燥完成。与冻结初始扩散系数的对照相比，前3 h含水率最大差为\textbf{0.3981 kg/kg}，说明扩散系数随状态变化会明显影响水分分布。

\textbf{针对问题三，}假设4 h后烘房维持50 $^\circ$C和0.05 kg/kg，以\textbf{全域最大干基含水率}严格低于0.15 kg/kg为达标判据，结合\textbf{事件监测、局部求根与严格达标复核}，得到烘干时长\textbf{57.4731 h}，中心控制最终达标。后期温度$\pm2\ ^\circ$C情景使临界时长变化为$-6.08\%$至$+6.65\%$。在端面与侧面交换系数相等的设定下，\textbf{二维轴对称配对校核}的临界时间比一维提前约\textbf{8 s（0.004\%）}，表明该情景下端面对达标时长的影响较小，支持径向简化的合理性。

\textbf{针对问题四，}在\textbf{比例径向收缩假设}下，结合给定半径轨迹确定材料速度场，再由干物质与水质量守恒建立\textbf{材料坐标模型}，将移动区域映射至固定区间。采用附录4物性，得到烘干时长\textbf{51.0877 h}、终时半径1.2000 cm。同物性、同环境下的固定半径对照时长为129.8448 h，收缩使模型预测时长缩短\textbf{78.7571 h（60.65\%）}，量化了给定收缩轨迹对干燥进程的影响。

参考解对照、网格与时间加密支持数值结果在已检验配置下的一致性，四问水分收支相对残差均小于$10^{-6}$。所报时长为给定输入与模型假设下的预测，尚需内部温湿观测验证，未显式计入的潜热及界面平衡关系仍构成适用限制。

\vspace{0.5em}
\noindent\textbf{关键词：}热湿耦合；变物性；有限体积法；全域达标；径向收缩
